\documentclass[11pt,a4paper]{article}

\usepackage[T1]{fontenc}
\usepackage[utf8]{inputenc}
\usepackage{lmodern}
\usepackage{amssymb,amsmath,amsthm,mathtools}
\usepackage[margin=2.5cm]{geometry}
\usepackage{parskip}
\usepackage[hidelinks]{hyperref}
\usepackage{microtype}

% Theorem environments (arXiv / journal style)
\theoremstyle{plain}
\newtheorem{theorem}{Theorem}[section]
\newtheorem{proposition}[theorem]{Proposition}
\newtheorem{lemma}[theorem]{Lemma}
\newtheorem{corollary}[theorem]{Corollary}

\theoremstyle{definition}
\newtheorem{definition}[theorem]{Definition}
\newtheorem{example}[theorem]{Example}
\newtheorem{remark}[theorem]{Remark}

% Notation
\newcommand{\Type}{\mathbf{Type}}
\newcommand{\id}{\mathrm{id}}
\newcommand{\Update}{\mathsf{Update}}

\title{Inventory as an Update Monad:\\
  Constant-Position Directed Containers}
\author{Extracted from a machine-checked development}
\date{}

\begin{document}
\maketitle

\begin{abstract}
Directed containers whose position sets are constant are exactly monoid
actions. We record the equivalence in both directions, embed the
constant-position presentation into the general theory of directed
containers, and cointerpret the resulting structure as an \emph{update
monad} $TX = S \to M \times X$ whose Kleisli maps observe a stock level
and emit a movement. Two supply-chain remarks are machine-checked:
abelianization of movement words forgets audit-trail order, and truncated
subtraction on~$\mathbb{N}$ clamps unmet demand (repaired by passing
to~$\mathbb{Z}$). A two-product warehouse with an order-up-to policy
illustrates the Kleisli reading. The proofs of all stated results have
been machine-checked; we record only the mathematical content.
\end{abstract}

\tableofcontents
\bigskip

%======================================================================
\section{Introduction}
%======================================================================

Inventory management tracks a stock level and a monoid of stock
movements. Applying a movement updates the stock; composing movements
associates; the empty movement does nothing. This is a right monoid
action of movements on stock levels.

The same data may be read as a directed container whose positions do not
depend on the shape: the single position set is the monoid of movements,
the root is the unit, and the shift is multiplication. Conversely, every
constant-position directed container is a monoid action. Under this
correspondence the cointerpretation of Ahman and Uustalu becomes the
\emph{update monad}
\[
  TX \;=\; S \longrightarrow M \times X,
\]
whose Kleisli arrows observe a state and emit an update together with a
result---precisely the form of an inventory policy that looks at current
stock and decides what to order.

\paragraph{What this note claims.}
In ordinary mathematical language:
\begin{enumerate}
  \item monoid actions $\longleftrightarrow$ constant-position directed
        containers, with both round trips definitional;
  \item the embedding of the constant-position case into general directed
        containers is transport-free once one uses that transport along a
        shape equality is the identity when positions are constant;
  \item the update monad of a monoid action satisfies the three monad
        laws;
  \item abelianization of free movement words forgets order (audit trail
        versus balance), and monus on~$\mathbb{N}$ clamps backorders
        while~$\mathbb{Z}$ does not.
\end{enumerate}

We work throughout in a single universe of types. Proofs are indicated by
the phrase ``by lean,'' meaning they have been checked by a proof
assistant rather than written out in prose.

%======================================================================
\section{Monoid actions}
%======================================================================

\begin{definition}[Monoid]\label{def:monoid}
A \textbf{monoid} $M$ consists of a type $|M|$, a multiplication
$\cdot : |M| \times |M| \to |M|$, and a unit $1 \in |M|$, satisfying
associativity and the two unit laws.
\end{definition}

\begin{definition}[Monoid action]\label{def:action}
A \textbf{right monoid action} consists of a monoid~$M$, a type~$S$ of
states, and a map
\[
  \mathbin{\triangleleft} \;:\; S \times |M| \;\longrightarrow\; S
\]
satisfying
\begin{equation}
  s \mathbin{\triangleleft} 1 \;=\; s,
  \qquad
  (s \mathbin{\triangleleft} m) \mathbin{\triangleleft} m'
  \;=\;
  s \mathbin{\triangleleft} (m \cdot m').
\end{equation}
For inventory, $S$ is the set of stock levels and $M$ the monoid of stock
movements; $s \mathbin{\triangleleft} m$ applies the movements~$m$ to the
stock~$s$.
\end{definition}

%======================================================================
\section{Constant-position directed containers}
%======================================================================

\begin{definition}[Constant-position directed container]\label{def:const-dc}
A \textbf{constant-position directed container} consists of
\begin{enumerate}
  \item a type $S$ of shapes (states),
  \item a single type $M$ of positions,
  \item a sub-shape map $\mathsf{sub} : S \to M \to S$,
  \item a root position $r \in M$ (independent of shape),
  \item a shift operation $\mathsf{shift} : M \to M \to M$
        (independent of shape),
\end{enumerate}
subject to the five laws, stated without transport:
\begin{align}
  \mathsf{sub}\,s\,r &\;=\; s, \label{eq:c1}\\
  \mathsf{shift}\,r\,m &\;=\; m, \label{eq:c2}\\
  \mathsf{shift}\,m\,r &\;=\; m, \label{eq:c3}\\
  \mathsf{sub}\,s\,(\mathsf{shift}\,p\,q)
    &\;=\; \mathsf{sub}\,(\mathsf{sub}\,s\,p)\,q, \label{eq:c4}\\
  \mathsf{shift}\,(\mathsf{shift}\,p\,q)\,r'
    &\;=\; \mathsf{shift}\,p\,(\mathsf{shift}\,q\,r'). \label{eq:c5}
\end{align}
\end{definition}

\begin{remark}[Why constant positions matter]
With constant positions, the shift is a genuine binary operation on the
single set~$M$. In the general directed-container case, positions vary
with shape, so shift is only a fibrewise composition and no global monoid
exists on positions. The constant-position case is the sound special
case of that observation.
\end{remark}

\begin{remark}
We deliberately do \emph{not} formulate the structure as a general
directed container plus a type equality $\forall s.\; \mathsf{Pos}(s)=M$.
That formulation buries every law in transports. The presentation above
keeps the equalities homogeneous.
\end{remark}

%======================================================================
\section{The equivalence}
%======================================================================

\begin{definition}[From constant-position DC to action]\label{def:to-action}
Given a constant-position directed container~$C$, define a monoid action
by taking $|M| = C.M$, multiplication $= C.\mathsf{shift}$, unit $=
C.r$, states $= C.S$, and action $= C.\mathsf{sub}$. Laws~\eqref{eq:c2},
\eqref{eq:c3}, \eqref{eq:c5} are the monoid laws;
laws~\eqref{eq:c1},~\eqref{eq:c4} are the action laws.
\end{definition}

\begin{definition}[From action to constant-position DC]\label{def:to-const}
Given a monoid action of~$M$ on~$S$, define a constant-position directed
container by taking shapes~$S$, positions~$|M|$, sub-shape $=$ the
action, root $=$ the monoid unit, and shift $=$ monoid multiplication.
\end{definition}

\begin{theorem}[Round trips]\label{thm:roundtrips}
The constructions of Definitions~\ref{def:to-action}
and~\ref{def:to-const} are mutually inverse:
\begin{equation}
  A \longmapsto \text{const} \longmapsto \text{action}
  \;=\; A,
  \qquad
  C \longmapsto \text{action} \longmapsto \text{const}
  \;=\; C.
\end{equation}
\end{theorem}

\begin{proof}
By lean.
\end{proof}

\begin{remark}
The equalities are definitional: both directions are dictionaries
renaming the same data, and the round trips hold by structure $\eta$ and
proof irrelevance.
\end{remark}

%======================================================================
\section{Embedding into directed containers}
%======================================================================

\begin{lemma}[Transport is trivial]\label{lem:transport-const}
Let $S$ and $M$ be types, and let $h : s = t$ be an equality of shapes.
For every $q \in M$,
\begin{equation}
  \mathsf{transport}_{h}(q) \;=\; q.
\end{equation}
That is, transport of a position along a shape equality is the identity
when the position family is constant.
\end{lemma}

\begin{proof}
By lean.
\end{proof}

\begin{definition}[Embedding]\label{def:to-directed}
Every constant-position directed container~$C$ determines a directed
container with
\begin{enumerate}
  \item shapes $C.S$,
  \item $\mathsf{Pos}(s) = C.M$ for every shape~$s$,
  \item root constantly equal to $C.r$,
  \item sub-shape and shift as in~$C$.
\end{enumerate}
\end{definition}

\begin{theorem}[Laws match]\label{thm:to-directed}
Under the embedding of Definition~\ref{def:to-directed}, the five
directed-container axioms are exactly the five constant-position laws of
Definition~\ref{def:const-dc}, with the D2/D5 transports erased by
Lemma~\ref{lem:transport-const}.
\end{theorem}

\begin{proof}
By lean.
\end{proof}

\begin{corollary}[Monoid actions as directed containers]\label{cor:action-dc}
Every monoid action determines a directed container, by composing
Definition~\ref{def:to-const} with Definition~\ref{def:to-directed}.
Directed containers whose position sets are constant are exactly monoid
actions.
\end{corollary}

\begin{proof}
By lean.
\end{proof}

%======================================================================
\section{The update monad}
%======================================================================

\begin{definition}[Update monad]\label{def:update}
Let $A$ be a monoid action of monoid~$M$ on states~$S$. The
\textbf{update monad} of~$A$ is the endofunctor
\begin{equation}
  \Update_A X \;=\; S \longrightarrow |M| \times X.
\end{equation}
A Kleisli map $X \to \Update_A Y$ observes a state in~$S$ and emits an
update in~$|M|$ together with a result in~$Y$.
\end{definition}

\begin{definition}[Unit and bind]\label{def:pure-bind}
\begin{enumerate}
  \item $\mathsf{pure} : X \to \Update_A X$ emits the monoid unit and
        returns the value:
        \[
          \mathsf{pure}(x)(s) \;=\; (1,\, x).
        \]
  \item Given $t : \Update_A X$ and $k : X \to \Update_A Y$,
        \[
          \mathsf{bind}(t,k)(s)
          \;=\;
          \bigl(
            m \cdot m',\;
            y
          \bigr),
        \]
        where $(m,x) = t(s)$, $(m',y) = k(x)\bigl(s \mathbin{\triangleleft} m\bigr)$:
        run~$t$, apply its update to the state, run~$k$ on the updated
        state, and multiply the two updates.
\end{enumerate}
\end{definition}

\begin{theorem}[Monad laws]\label{thm:monad}
For every monoid action~$A$ and types $X,Y,Z$:
\begin{enumerate}
  \item \textbf{Left unit:}
        $\mathsf{bind}(\mathsf{pure}(x),\, k) = k(x)$;
  \item \textbf{Right unit:}
        $\mathsf{bind}(t,\, \mathsf{pure}) = t$;
  \item \textbf{Associativity:}
        $\mathsf{bind}(\mathsf{bind}(t,k),\, h)
         = \mathsf{bind}\bigl(t,\, x \mapsto \mathsf{bind}(k(x),h)\bigr)$.
\end{enumerate}
\end{theorem}

\begin{proof}
By lean.
\end{proof}

\begin{remark}
The left unit uses the action unit law (the state is unchanged) and the
monoid left unit law. The right unit uses the monoid right unit law.
Associativity uses the action multiplication law (the middle state is the
same under either bracketing) and monoid associativity.
\end{remark}

\begin{remark}[Cointerpretation]
The update monad is the cointerpretation of the constant-position
directed container associated to~$A$, in the sense of Ahman and Uustalu:
directed containers cointerpret as monads, and the constant-position case
specialises to update monads.
\end{remark}

%======================================================================
\section{Two remarks on movements and stock}
%======================================================================

\begin{definition}[Movement words]\label{def:movements}
A \textbf{stock movement} for a single product is either a receipt of
quantity $q \in \mathbb{N}$ or an issue of quantity~$q$. Words in the free
monoid on movements are the storekeeper's \textbf{audit trail}. The
\textbf{net effect} of a word is its image in the abelianization
$\mathbb{Z}$ (the ledger \textbf{balance}):
\begin{equation}
  \mathsf{net}([]) = 0,\quad
  \mathsf{net}(\mathsf{receipt}\,q :: w) = q + \mathsf{net}(w),\quad
  \mathsf{net}(\mathsf{issue}\,q :: w) = \mathsf{net}(w) - q.
\end{equation}
\end{definition}

\begin{proposition}[Abelianization forgets order]\label{prop:abelian}
Receive-$2$-then-issue-$1$ and issue-$1$-then-receive-$2$ have the same
net balance but are distinct as lists:
\begin{equation}
  \mathsf{net}\bigl([\mathsf{receipt}\,2,\, \mathsf{issue}\,1]\bigr)
  \;=\;
  \mathsf{net}\bigl([\mathsf{issue}\,1,\, \mathsf{receipt}\,2]\bigr),
\end{equation}
while the two words are unequal. Operationally the difference matters:
from a stock of zero the second word passes through a shortage the first
never sees.
\end{proposition}

\begin{proof}
By lean.
\end{proof}

\begin{proposition}[Monus clamps; integers admit backorders]\label{prop:monus}
\begin{enumerate}
  \item There exist $s,m \in \mathbb{N}$ with
        $(s \mathbin{\dot{-}} m) + m \neq s$ (truncated subtraction
        silently discards unmet demand).
  \item For all $s,m \in \mathbb{Z}$,
        $(s - m) + m = s$ (negative stock is a backorder; every issue is
        exactly undone by the matching receipt).
\end{enumerate}
\end{proposition}

\begin{proof}
By lean.
\end{proof}

%======================================================================
\section{Example: a two-product warehouse}
%======================================================================

\begin{example}[Inventory action]\label{ex:warehouse}
Let $M = \mathbb{Z} \times \mathbb{Z}$ under componentwise addition (net
movements for two products; the abelianization $\mathbb{Z}^{\mathrm{Prod}}$
with $|\mathrm{Prod}| = 2$), and let $S = \mathbb{Z} \times \mathbb{Z}$
(stock levels, backorders admitted). The action is componentwise
addition:
\begin{equation}
  (s_1,s_2) \mathbin{\triangleleft} (m_1,m_2)
  \;=\;
  (s_1+m_1,\, s_2+m_2).
\end{equation}
The associated directed container has stock levels as shapes and
movements as positions: a receipt followed by an equal issue is a
nontrivial endomorphism of the stock state.
\end{example}

\begin{definition}[Order-up-to policy]\label{def:policy}
The Kleisli map $\mathsf{policy} : 1 \to \Update_{\mathrm{inv}} 1$ is the
replenishment rule: for each product, if stock is below the reorder point
$5$, order up to~$10$; otherwise order nothing. Goods inward
$\mathsf{receive}(m)$ is the Kleisli map that unconditionally emits the
movement~$m$.
\end{definition}

\begin{proposition}[Policy behaviour]\label{prop:policy}
\begin{enumerate}
  \item At stock $(2,7)$, the policy emits movement $(8,0)$.
  \item Applying that movement yields stock $(10,7)$.
  \item Kleisli composition threads state: receiving $(3,0)$ before the
        policy, starting from stock $(1,7)$, produces total movement
        $(9,0)$ (the policy sees stock $4$ and orders~$6$).
\end{enumerate}
\end{proposition}

\begin{proof}
By lean.
\end{proof}

%======================================================================
\section{Discussion}
%======================================================================

\paragraph{What has been established.}
Constant-position directed containers and monoid actions are the same
data. Their cointerpretation is the update monad, with the three monad
laws following from the monoid and action laws. The two remarks separate
the free monoid of audit trails from its abelianization (balances), and
motivate working over~$\mathbb{Z}$ rather than~$\mathbb{N}$ when unmet
demand must be recorded.

\paragraph{Scope.}
The development treats a single universe of types and does not address
multi-echelon inventory, stochastic demand, or cost optimisation. The
warehouse example is a pure illustration of the Kleisli reading.

\paragraph{Machine checking.}
Every definition and theorem statement above corresponds to a declaration
in a machine-checked formalisation. The phrase ``by lean'' indicates that
the proof has been verified by a proof assistant.

%======================================================================
\section*{Acknowledgements}
%======================================================================

The mathematical content follows the cointerpretation of directed
containers as monads due to Ahman and Uustalu, specialised to the
constant-position case and read as inventory. The presentation here is a
prose rendering of a formal development intended for human review of the
theorem statements.

\begin{thebibliography}{9}

\bibitem{ahman2016}
D.~Ahman and T.~Uustalu.
\newblock Directed containers as categories.
\newblock In \emph{MSFP}, volume~207 of \emph{EPTCS}, pages 89--98, 2016.

\bibitem{ahman-uustalu-update}
D.~Ahman and T.~Uustalu.
\newblock Update monads: cointerpreting directed containers.
\newblock In \emph{TLCA}, volume~8560 of \emph{LNCS}, pages 1--15.
\newblock Springer, 2014.

\bibitem{abbott2005}
M.~Abbott, T.~Altenkirch, and N.~Ghani.
\newblock Containers: constructing strictly positive types.
\newblock \emph{Theoretical Computer Science}, 342(1):3--27, 2005.

\end{thebibliography}

\end{document}
